\chapter[Band 43: Ringe mit Eins]{Band 43 auf einen Blick}
\noindent\textbf{Ringe mit Eins}\par\medskip
Ein Ring verbindet die abelschen Gruppen aus Band~41 mit Monoiden.
Die Multiplikation muss nicht kommutativ sein; \(0\ne1\) wird nicht vorausgesetzt.

\begin{BandUebersicht}
\textbf{Ringaxiome}

\UebFormel{\begin{aligned}&\AbelianGroupAlg R+0,\quad\MonoidAlg R\cdot1,\\&x(y+z)=xy+xz,\\&(x+y)z=xz+yz.\end{aligned}}
\UebQuellen{\FormulaRefAuto{RingDef}[def]}
&
\textbf{Null und Vorzeichen}

\UebFormel{\begin{aligned}&a0=0=0a,\quad-(-a)=a,\\&(-a)b=-(ab)=a(-b),\\&(-a)(-b)=ab.\end{aligned}}
\UebQuellen{\FormulaRefAuto{RingLeftZeroAbsorption}[thm];
\FormulaRefAuto{RingRightZeroAbsorption}[thm];
\FormulaRefAuto{RingDoubleNegative}[thm];
\FormulaRefAuto{RingNegativeTimesNegative}[thm]} \\
\addlinespace[.8ex]

\textbf{Ringhomomorphismen}

\UebFormel{\begin{aligned}&f(a+b)=f(a)+f(b),\\&f(ab)=f(a)f(b),\quad f(1_R)=1_S.\end{aligned}}
\UebQuellen{\FormulaRefAuto{RingHomDef}[def];
\FormulaRefAuto{RingIsoDef}[def]}
&
\textbf{Folgerungen}
Kompositionen sind wieder Ringhomomorphismen.
\UebFormel{f(0_R)=0_S,\qquad f(-a)=-f(a).}
\UebQuellen{\FormulaRefAuto{RingHomZeroPreservation}[thm];
\FormulaRefAuto{RingHomNegativePreservation}[thm];
\FormulaRefAuto{RingHomComposition}[thm]} \\
\addlinespace[.8ex]

\textbf{Ganzzahliger Ring}

\UebFormel{\RingAlg{\mathbb Z}{+}{\cdot}{\IntZero}{\IntOne}.}
\UebQuellen{\FormulaRefAuto{IntegerUnitalRing}[thm]}
&
\textbf{Kanonische Abbildung}
\(\eta_R\) ist die natürliche Halbringabbildung aus Band~39.
\UebFormel{\Phi_R(\IntClass a b)=\eta_R(a)\oplus(-\eta_R(b)).}
\UebQuellen{\FormulaRefAuto{IntegerRingMapDef}[def];
\FormulaRefAuto{IntegerRingMapClassValue}[thm]} \\
\addlinespace[.8ex]

\textbf{Erhaltung der Operationen}

\UebFormel{\begin{aligned}&\Phi_R(z+w)=\Phi_R(z)\oplus\Phi_R(w),\\&\Phi_R(zw)=\Phi_R(z)\odot\Phi_R(w).\end{aligned}}
\UebQuellen{\FormulaRefAuto{IntegerRingMapAddition}[thm];
\FormulaRefAuto{IntegerRingMapMultiplication}[thm];
\FormulaRefAuto{IntegerRingMapConstants}[thm]}
&
\textbf{Universelle Eigenschaft}
Für jeden Ring mit Eins ist \(\Phi_R\) der eindeutig bestimmte Ringhomomorphismus aus \(\mathbb Z\).
\UebFormel{\exists!f\,\RingHom f{\mathbb Z,+,\cdot,\IntZero,\IntOne}{R,\oplus,\odot,0_R,1_R}.}
\UebQuellen{\FormulaRefAuto{IntegerRingUniversalHomomorphism}[thm]} \\
\addlinespace[.8ex]
\end{BandUebersicht}
